\section{Mathematical Model \& Cryptographic Primitives}
\label{sec:methodology}

\subsection{HKDF Sub-Key Derivation Pipeline}
\label{subsec:hkdf_pipeline}
The key schedule subsystem derives three domain-separated sub-keys from a master key $K$, a 16-byte random salt $S$, and a 12-byte nonce $N$:
\begin{align}
\text{PRK} &= \text{HMAC-SHA256}(S, K) \label{eq:hkdf_extract} \\
K_r &= \text{HKDF-Expand}(\text{PRK}, \text{"KDR-CA-AEAD-v1-ca-rules|"} \,\|\, N, 32) \label{eq:hkdf_kr} \\
K_c &= \text{HKDF-Expand}(\text{PRK}, \text{"KDR-CA-AEAD-v1-cipher-key|"} \,\|\, N, 32) \label{eq:hkdf_kc} \\
K_a &= \text{HKDF-Expand}(\text{PRK}, \text{"KDR-CA-AEAD-v1-mac-key|"} \,\|\, N, 32) \label{eq:hkdf_ka}
\end{align}

Here, $K_r \in \{0,1\}^{256}$ is formatted into a 32-element tuple of uint8 integers $R = (r_0, r_1, \dots, r_{31})$, where each $r_j \in [0, 255]$. The 32 rules directly correspond to the 256 bits of key material derived from HKDF-SHA256.

\subsection{1D Elementary Cellular Automaton Evaluation}
\label{subsec:eca_eval}
For each byte position index $i \in [0, M-1]$ in the payload stream of length $M$:
\begin{align}
R_1 &= R[i \bmod 32] \label{eq:r1_select} \\
R_2 &= R[(i + \delta) \bmod 32] \label{eq:r2_select}
\end{align}
where $\delta = 13$ is the prime offset index ensuring non-aligned dual-rule coupling.

The local ECA byte evaluation $S_{\text{ECA}}(i, R_1, R_2)$ computes bit-level Wolfram neighborhood transitions. For each bit position $b \in [0, 7]$ within the byte:
\begin{equation}
\text{state}_{\text{init}} = (i \oplus R_1) \bmod 256
\end{equation}
The left, center, and right neighborhood bits of $\text{state}_{\text{init}}$ at bit position $b$ under periodic boundary conditions are:
\begin{align}
\text{left} &= (\text{state}_{\text{init}} \gg ((b+1) \bmod 8)) \ \& \ 1 \\
\text{self} &= (\text{state}_{\text{init}} \gg b) \ \& \ 1 \\
\text{right} &= (\text{state}_{\text{init}} \gg ((b-1) \bmod 8)) \ \& \ 1
\end{align}
Neighborhood index:
\begin{equation}
\eta = (\text{left} \ll 2) \ |\  (\text{self} \ll 1) \ |\  \text{right} \in [0, 7]
\end{equation}
The output bit $b_{\text{out}}$ is selected from secondary rule $R_2$:
\begin{equation}
b_{\text{out}} = (R_2 \gg \eta) \ \& \ 1
\end{equation}
Reassembling all 8 output bits yields $S_{\text{ECA}} = \left( \sum_{b=0}^7 b_{\text{out}} \cdot 2^b \right) \oplus R_1$.

\subsection{Candidate A-Chain State Permutation}
\label{subsec:candidate_a_chain}
Given input plaintext byte $P_i$ and previous feedback state $\text{prev\_state}$ (initialized to $\text{IV}_0 = \text{0xC5}$):

\begin{enumerate}
    \item \textbf{Inter-Byte State Chaining & Modulo Addition}:
    \begin{equation}
    y_1 = ((P_i \oplus \text{prev\_state}) + S_{\text{ECA}}) \bmod 256
    \label{eq:chain_mod}
    \end{equation}
    
    \item \textbf{Keyed Circular Right Rotation}:
    \begin{equation}
    \alpha = (R_1 \bmod 7) + 1 \in [1, 7]
    \end{equation}
    \begin{equation}
    y_2 = \text{ROTR}_8(y_1, \alpha) = ((y_1 \gg \alpha) \ |\  (y_1 \ll (8 - \alpha))) \ \& \ \text{0xFF}
    \label{eq:keyed_rotr}
    \end{equation}
    
    \item \textbf{Secondary Rule Mixing & Feedback Update}:
    \begin{equation}
    T_i = y_2 \oplus R_2
    \label{eq:sec_rule_mix}
    \end{equation}
    \begin{equation}
    \text{prev\_state} \leftarrow T_i
    \end{equation}
\end{enumerate}

\subsection{Reversible Inverse Permutation Pipeline}
\label{subsec:inverse_permutation}
Decryption applies the exact inverse transformations in reverse order:
\begin{align}
y_2 &= T_i \oplus R_2 \label{eq:inv_rule_mix} \\
y_1 &= \text{ROTL}_8(y_2, \alpha) = ((y_2 \ll \alpha) \ |\  (y_2 \gg (8 - \alpha))) \ \& \ \text{0xFF} \label{eq:inv_rotr} \\
\text{mixed\_byte} &= (y_1 - S_{\text{ECA}}) \bmod 256 \label{eq:inv_mod} \\
P_i &= \text{mixed\_byte} \oplus \text{prev\_state} \label{eq:inv_chain} \\
\text{prev\_state} &\leftarrow T_i \label{eq:inv_feedback_update}
\end{align}

Because every step is bijective over $\mathbb{Z}_{256}$, exact round-trip reconstruction $P_i = \text{Inverse}(\text{Forward}(P_i))$ is mathematically guaranteed.
